%% ===== Discussion (~500-650 words) =====
\section{Discussion\cb{target ~570 words}\wc{discussion}\secstatus{todo}}\label{sec:discussion}
\textit{[FILLER -- Discussion.]} A strain-aware, biologically grounded
representation is the differentiator; limitations include single-organism scope,
morphology-data maturity, and error-bar provenance.

A central modeling assumption deserves explicit statement. TorchCell treats the catalogued
biological networks (Fig.~\ref{fig:torchcell}c) as annotations of a fixed reference and assumes
the wiring is perturbation-invariant, so a strain is a thin operator applied to the reference
rather than a rebuilt graph. This is a hypothesis, not a fact: an edge can effectively vanish
when, for example, perturbing an upstream regulator drives a partner's expression so low that
the interaction no longer occurs in the cell, and such condition-specific rewiring is not
captured by a static adjacency. We adopt the assumption deliberately---it makes the
perturbation a learnable displacement in representation space, keeps the biological graphs as a
soft prior the data can overrule rather than a hard substrate, and yields a single general data
object that different modeling paradigms can consume---and we treat its biological limits,
computational advantages, and pipeline generality at length in \suppnoteref{note:static-graph}.
